% Annotatsiooni tekst läheb siia (square brackets should be removed)
Euroopa liidu eFTI (Electronic Freight Transport Information) regulatsioon näeb ette veoinfo digitaalse
vahetamise pädevate asutuste ja veoettevõtete vahel. Liikmesriigid peavad olema aastaks 2027 valmis vastu võtma
veoinfot digitaalsel kujul ning muutma selle kättesaadavaks pädevatele asutustele.
Sellest tulenevalt peab olema igal liikmesriigil oma eFTI värav, mis infot osapoolte vahel edastama hakkab.

Regulatsiooni kohaselt on liikmesriikide ja riigisisese andmevahetuse aluseks eDelivery infovahetuse süsteem.
Riikide vahel kasutatakse staatilist eDelivery konfiguratsiooni,
mille kohaselt seadistatakse vajaminevad sertifikaadid ning muud parameetrid manuaalselt.
eFTI värav peab võimaldama veoettevõtetele ehk eFTI platvormidele ka eDelivery liidest,
kuid võib pakkuda ka teisi lahendusi nagu näiteks REST liides.

Platvormidele loodud eDelivery lahendus peab kasutama dünaamilist konfiguratsiooni, mis automatiseerib
sertifikaatide ning muude parameetrite leidmist.
Dünaamiliselt metaandmete leidmiseks on eDeliverys kasutusel kaks lisakomponenti: SMP (Teenuse metaandmete avaldaja) ja SML (Teenuse metaandmete leidja).

Käesolev lõputöö uurib, kuidas eDelivery dünaamiline otsing praktikas toimib, ning realiseerib töötava prototüübi.
Töö analüüsib olemasolevaid avatud lähtekoodiga komponente, toob välja praktilised integratsioonipiirangud
ning implementeerib minimaalse Kotlinil põhineva SMP/SML lahenduse, mis toetab Harmony Access Pointi
dünaamilise otsingu nõudeid.

Valminud prototüüp demonstreerib kahe juurdepääsupunkti vahelist täielikku dünaamilist sõnumivahetust ning seda hinnatakse
funktsionaalsete testide ja jõudlusmõõtmiste abil. Tulemused näitavad, et fokusseeritud ja minimaalsete
sõltuvustega teostus suudab lokaalses testkeskkonnas pakkuda suurt läbilaskevõimet nii SMP metaandmete vastustele
kui ka SML DNS-päringutele.



% No need to change this
Lõputöö on kirjutatud \langEst~keeles ning sisaldab teksti \calculatepages leheküljel, 
\total{totalchapters} peatükki\ifthenelse{\equal{\totvalue{figure}}{0}}{}{% If no figures, do nothing
, \total{figure} \ifnum\totvalue{figure}=1 joonis\else joonist\fi%
}\ifthenelse{\equal{\totvalue{table}}{0}}{}{% If no tables, do nothing
, \total{table} \ifnum\totvalue{table}=1 tabel\else tabelit\fi%
}.
